\chapter{Diskussion}\label{diskussion}

\begin{center}\rule{0.5\linewidth}{0.5pt}\end{center}

\section{Einordnung in den Stand der
Technik}\label{einordnung-in-den-stand-der-technik}

\subsection{7.1.1 Abgrenzung zu bestehenden
Systemen}\label{abgrenzung-zu-bestehenden-systemen}

Das in dieser Arbeit entwickelte System unterscheidet sich von
bestehenden Ansätzen der automatisierten Sportberichterstattung (Kapitel
3.3) in mehreren Aspekten:

\textbf{Gegenüber templatebasierten Systemen} (z. B. Retresco, vgl.
Beils 2023, S. 207) bietet der LLM-basierte Ansatz eine höhere
sprachliche Varianz und Anpassungsfähigkeit an unterschiedliche
Stilprofile. Templatebasierte Systeme erzeugen deterministische, aber
oft formelhaft wirkende Texte; LLM-generierte Texte sind sprachlich
natürlicher, dafür nicht-deterministisch.

\textbf{Gegenüber vollautonomen KI-Systemen} betont die vorliegende
Arbeit bewusst den hybriden Ansatz mit Human-in-the-Loop-Kontrolle. Dies
begründet sich durch die in Kapitel 3.1 beschriebene
Halluzinationsproblematik, die im journalistischen Kontext besonders
kritisch ist (Bluhm \& Schäfer 2023, S. 33). Der \texttt{coop}-Modus
repräsentiert damit einen Mittelweg, der die Effizienzgewinne der
Automatisierung nutzt, ohne die publizistische Verantwortung an die KI
zu delegieren.

\textbf{Gegenüber reinen Übersetzungslösungen} für Mehrsprachigkeit
bietet das System eine nativ multilinguale Generierung, bei der die
Zielsprache als Prompt-Parameter übergeben wird. Die zusätzliche
Batch-Übersetzungsfunktion mit reduzierter Temperatur (0,1) ergänzt dies
für nachträgliche Lokalisierung bestehender Einträge.

\subsection{7.1.2 Beitrag zur Forschung}\label{beitrag-zur-forschung}

Die Arbeit leistet drei Beiträge zum aktuellen Forschungsstand:

\begin{enumerate}
\def\labelenumi{\arabic{enumi}.}
\item
  \textbf{Architektonischer Beitrag}: Das System demonstriert eine
  produktionsnahe Referenzarchitektur für KI-gestützte
  Echtzeit-Textgenerierung mit modularer Provider-Abstraktion,
  Few-Shot-Prompting aus einer kuratierten Stilreferenz-Datenbank und
  dreistufigem Betriebsmodus-Konzept.
\item
  \textbf{Methodischer Beitrag}: Die Evaluationsinfrastruktur mit
  Bulk-Generierungsendpunkt, Provider-Override und statistischen
  Vergleichsmetriken (Cliff's Delta, Bootstrap-CI, Cohen's Kappa)
  ermöglicht systematische, reproduzierbare Modellvergleiche.
\item
  \textbf{Praktischer Beitrag}: Die White-Label-Architektur mit
  instanzspezifischer Stilsteuerung adressiert den in Kapitel 2.3
  beschriebenen Bedarf von Vereinen als eigenständige Medienproduzenten.
\end{enumerate}

\begin{center}\rule{0.5\linewidth}{0.5pt}\end{center}

\section{Kritische Reflexion}\label{kritische-reflexion}

Die in Kapitel 6.7 dokumentierten Limitationen werden im Folgenden
hinsichtlich ihrer Implikationen eingeordnet.

\subsection{7.2.1 Methodische Einordnung}\label{methodische-einordnung}

Die zentrale methodische Einschränkung betrifft die Selbstbewertung der
Textqualität (vgl. Kapitel 6.7, „Keine externe Nutzerstudie''). Im
Design Science Research ist die Evaluation durch den Entwickler ein
verbreitetes Vorgehen bei initialen Artefaktiterationen (Hevner et
al.~2004), begrenzt jedoch die externe Validität der qualitativen
Ergebnisse. Die implementierte Cohen's-Kappa-Metrik ist bewusst als
Infrastruktur für eine Folgestudie mit externen Ratern angelegt.

Die eingeschränkte Stichprobengröße (vgl. Kapitel 6.7, „Eingeschränkte
Stichprobengröße'') relativiert die Generalisierbarkeit der
quantitativen Ergebnisse. Insbesondere Randsituationen --- torloses
Unentschieden, Elfmeterschießen, Spielabbrüche --- sind in der
Datenbasis vermutlich unterrepräsentiert. Die
Bulk-Evaluationsinfrastruktur ermöglicht jedoch eine systematische
Erweiterung der Datenbasis ohne Codeänderung.

\subsection{7.2.2 Technische Einordnung}\label{technische-einordnung}

Das Polling-Modell (vgl. Kapitel 6.7, „Polling statt Push'') stellt für
den Einzelbetrieb keine Einschränkung dar, würde aber bei steigender
Nutzerzahl (\textgreater50 gleichzeitige Clients) zu einer relevanten
Serverlast führen. Die bestehende WebSocket-Infrastruktur für den
Media-Kanal zeigt, dass eine Umstellung auf Server-Sent Events für den
Ticker-Datenfluss technisch machbar ist.

Die fehlende Authentifizierung (vgl. Kapitel 6.7, „Keine
Authentifizierung'') ist im Kontext eines internen Redaktionswerkzeugs
vertretbar, schließt aber einen offenen Mehrbenutzerbetrieb ohne
Netzwerkabsicherung aus. Die Architektur ist durch die saubere
Middleware-Schichtung (FastAPI Dependencies) für eine nachträgliche
JWT-Integration vorbereitet.

Die 15 n8n-Workflows (vgl. Kapitel 6.7, „Keine Tests für
n8n-Workflows'') stellen den größten nicht automatisiert getesteten
Systemteil dar. Da sie die gesamte Datenversorgung des Systems
verantworten, wäre eine Absicherung durch Integrationstests gegen eine
Staging-Datenbank für den Produktivbetrieb prioritär.

Die Provider-Auswahl im LLM-Service ist als schlüsselbasierte
Priorisierung implementiert, nicht als konfigurierbare
Laufzeit-Strategie. Für einen Produktivbetrieb mit mehreren
gleichzeitigen Spielen wäre ein expliziter Konfigurationsmechanismus pro
Spiel (z. B. ein Datenbankfeld) robuster.

Der \texttt{useAutoPublisher}-Hook im Frontend enthält eine semantische
Doppelung: Im Auto-Modus setzt n8n über \texttt{auto\_publish=True} den
Status bereits beim Generieren; der Hook fängt Residualfälle auf. Diese
Zweispurigkeit entstand durch inkrementelle Entwicklung und könnte auf
eine definitive Strategie konsolidiert werden.

Die WebSocket-Verbindung existiert als In-Memory-Singleton pro Prozess.
Bei einem Multi-Prozess-Deployment (mehrere Uvicorn-Worker hinter einem
Load-Balancer) würden Medien-Updates nur Clients des jeweiligen Workers
erreichen; eine Redis-Pub/Sub-Erweiterung wäre die natürliche
Skalierungsstufe.

Der \texttt{strict}-Modus in \texttt{tsconfig.json} ist derzeit
deaktiviert, um die inkrementelle TypeScript-Migration nicht zu
blockieren. Eine schrittweise Aktivierung von \texttt{strictNullChecks}
und \texttt{noImplicitAny} ist der naheliegende Folgeschritt zur
weiteren Härtung der Codebasis.

\subsection{7.2.3 Inhaltliche Reflexion}\label{inhaltliche-reflexion}

Die in Kapitel 2.5 hergeleiteten linguistischen Anforderungen an
Liveticker --- Ellipsen, konzeptionelle Mündlichkeit, Graphostilistik
(z. B. „TOOOOR!{\kern0pt}``) --- stellen besondere Anforderungen an die
Prompt-Gestaltung. Die Erfahrung zeigt, dass LLMs dazu neigen, in einem
formelleren Register zu schreiben als es das Genre Liveticker erfordert.
Die Few-Shot-Referenzen aus der \texttt{style\_references}-Tabelle sind
das primäre Mittel, um diese stilistische Lücke zu schließen.

Die Evaluationsergebnisse (Abschnitt 6.3.6) zeigen, dass die
Few-Shot-Referenzen das Format zuverlässig konditionieren ---
Minutenangaben, TOOOOR-Konvention und Textkürze werden konsistent
übernommen. Die stilistische Lücke ist damit teilweise geschlossen:
Faktentreue und Verständlichkeit profitieren von der strukturierten
Kontextübergabe, während Tonalität als schwächste Dimension verbleibt.
Die häufigste Fehlerklasse --- Stil-Inkonsistenz (vgl.
Fehlerklassen-Tabelle in Abschnitt 6.3.6) --- tritt gerade dort auf, wo
euphorische Few-Shot-Muster in den neutralen Stil bluten. Eine Trennung
der Referenz-Pools nach Stilprofil wäre die naheliegende Korrektur.

Der in Kapitel 3.1 beschriebene Halluzinationseffekt ist im Kontext von
Livetickern besonders kritisch, da fehlerhafte Fakten (falscher
Torschütze, falsches Ergebnis) unmittelbar die Glaubwürdigkeit
zerstören. Die explizite Schutzregel für Pre-Match-Prompts und die
niedrige Temperatur (vgl. Abschnitt 7.4.2) sind Gegenmaßnahmen, deren
Wirksamkeit jedoch nur im \texttt{coop}-Modus durch die redaktionelle
Kontrolle vollständig abgesichert ist. Im \texttt{auto}-Modus verbleibt
ein Restrisiko fehlerhafter Veröffentlichungen.

\begin{center}\rule{0.5\linewidth}{0.5pt}\end{center}

\section{Diskussion der Betriebsmodi}\label{diskussion-der-betriebsmodi}

Die drei Betriebsmodi (\texttt{auto}, \texttt{coop}, \texttt{manual})
wurden in Kapitel 4.3.3 konzipiert und in Kapitel 6.4 evaluiert. Im
Folgenden werden die Implikationen der Evaluationsergebnisse für den
praktischen Einsatz diskutiert.

\subsection{7.3.1 Auto-Modus: Geschwindigkeit auf Kosten der
Kontrolle}\label{auto-modus-geschwindigkeit-auf-kosten-der-kontrolle}

Die Stärke des \texttt{auto}-Modus liegt in der Geschwindigkeit (vgl.
TTP-Messung in Abschnitt 6.4.2). Das Risiko besteht in unkontrollierten
Halluzinationen, die ohne redaktionelle Prüfung veröffentlicht werden.
Im journalistischen Kontext, in dem Glaubwürdigkeit eine zentrale
Ressource darstellt (Beils 2023, S. 57), ist dieser Modus daher nur für
unkritische Event-Typen (z. B. Phasenwechsel wie „Anpfiff'' oder
„Halbzeit'') vertretbar.

\subsection{7.3.2 Coop-Modus: Der intendierte
Produktivbetrieb}\label{coop-modus-der-intendierte-produktivbetrieb}

Im Evaluationszeitraum wurden im Coop-Modus keine Einträge nach Freigabe
zurückgezogen --- ein Hinweis darauf, dass die redaktionelle Prüfung als
Qualitätsstufe wirksam ist. Ob einzelne Entwürfe vor der Freigabe
inhaltlich editiert wurden, wurde nicht systematisch erfasst;
Bearbeitungen sind über das Frontend jedoch jederzeit möglich. Die
TAB/ESC-Tastatursteuerung erweist sich als zentrales Mittel, um den
Entscheidungsdurchsatz zu sichern, ohne den Redakteur visuell zu
überlasten.

\subsection{7.3.3 Manual-Modus: Status quo als
Vergleichsbasis}\label{manual-modus-status-quo-als-vergleichsbasis}

Der \texttt{manual}-Modus dient als Vergleichsbasis und Rückfalloption
--- etwa bei hochkritischen Ereignissen oder bei Redakteuren, die
zunächst Vertrauen in das System aufbauen möchten. Die
Laufzeit-Umschaltbarkeit per API-Call stellt sicher, dass der Wechsel
ohne Systemunterbrechung möglich ist.

\subsection{7.3.4 Implikationen für die
Praxis}\label{implikationen-fuxfcr-die-praxis}

Die Koexistenz aller drei Modi in einem System --- umschaltbar zur
Laufzeit per API-Call --- ermöglicht eine schrittweise Einführung in
Redaktionen: Neue Nutzer können im \texttt{manual}-Modus beginnen, über
den \texttt{coop}-Modus Vertrauen in die KI-Qualität aufbauen und für
unkritische Standardereignisse schließlich den \texttt{auto}-Modus
aktivieren. Diese Graduierung adressiert den in Kapitel 2.1
beschriebenen Akzeptanzbedarf bei der Einführung KI-gestützter Werkzeuge
in redaktionelle Workflows.

\begin{center}\rule{0.5\linewidth}{0.5pt}\end{center}

\section{Diskussion der
Prompt-Architektur}\label{diskussion-der-prompt-architektur}

\subsection{7.4.1 Few-Shot-Prompting als
Stilsteuerung}\label{few-shot-prompting-als-stilsteuerung}

Die dynamische Einbindung von bis zu drei Stilreferenzen aus der
\texttt{style\_references}-Datenbanktabelle stellt einen zentralen
Designentscheid dar. Im Vergleich zu einem statischen System-Prompt
bietet dieser Ansatz zwei Vorteile:

\begin{enumerate}
\def\labelenumi{\arabic{enumi}.}
\tightlist
\item
  \textbf{Anpassbarkeit ohne Codeänderung}: Neue Stilbeispiele können
  über die Datenbank hinzugefügt werden, ohne den Prompt-Code zu
  modifizieren.
\item
  \textbf{Instanzspezifik}: Die Filterung nach Event-Typ und Instanz
  (\texttt{ef\_whitelabel} vs.~\texttt{generic}) ermöglicht
  unterschiedliche Stilprofile für verschiedene Einsatzkontexte.
\end{enumerate}

Ein kontrollierter A/B-Test (mit vs.~ohne Referenzen) war nicht
durchführbar (vgl. Abschnitt 6.3.5). Qualitativ zeigt die Evaluation
konsistente Formatierungsmuster in der \texttt{ef\_whitelabel}-Instanz
--- insbesondere bei TOOOOR-Konvention und Minutenformat. Der Effekt ist
sichtbar, aber statistisch nicht isoliert messbar.

\subsection{7.4.2 Temperatur und
Determinismus}\label{temperatur-und-determinismus}

Die gewählte Temperatur von 0,3 für die Textgenerierung und 0,1 für
Übersetzungen ist ein Kompromiss zwischen sprachlicher Varianz und
Vorhersagbarkeit. Eine niedrigere Temperatur würde die
Halluzinationsrate weiter senken, aber die Texte formelhafter machen ---
ein Zielkonflikt, der im Genre Liveticker (das von Überraschung und
sprachlicher Kreativität lebt) besonders relevant ist.

\subsection{7.4.3 Kontext-Aufbereitung}\label{kontext-aufbereitung}

Die sieben spezialisierten Context-Builder (vgl. Abschnitt 5.2.6)
strukturieren die Fakten vor der Übergabe an das LLM. Diese
Vorverarbeitung reduziert die Wahrscheinlichkeit von Halluzinationen, da
das Modell nicht aus unstrukturierten Rohdaten extrahieren muss, sondern
bereits aufbereitete Faktenblöcke erhält. Die Wirksamkeit dieses
Ansatzes zeigt sich insbesondere bei Pre-Match-Einträgen, wo die
Faktengrundlage (z. B. Verletzungslisten, H2H-Statistiken) klar
abgegrenzt ist.

\begin{center}\rule{0.5\linewidth}{0.5pt}\end{center}

\section{Implikationen für den
Sportjournalismus}\label{implikationen-fuxfcr-den-sportjournalismus}

\subsection{7.5.1 Veränderung der
Redakteursrolle}\label{veruxe4nderung-der-redakteursrolle}

Das hybride System verschiebt die Rolle des Liveticker-Redakteurs von
der \textbf{Textproduktion} zur \textbf{Textredaktion}: Statt unter
Zeitdruck zu formulieren, prüft, editiert und kuratiert der Redakteur
KI-generierte Vorschläge. Diese Verschiebung entspricht dem von Bluhm
und Schäfer (2023, S. 35) beschriebenen Wandel hin zu einer „augmented
authorship'', bei der menschliche Expertise und maschinelle Effizienz
komplementär zusammenwirken.

Praktisch bedeutet dies, dass der kognitive Engpass der
Liveticker-Produktion --- die simultane Beobachtung des Spielgeschehens,
Relevanzbewertung und Texterstellung unter extremem Zeitdruck (vgl.
Kapitel 2.1) --- durch das System aufgelöst wird. Die originäre
Kompetenz des Redakteurs verlagert sich: Sprachliche
Formulierungsfähigkeit tritt in den Hintergrund, während evaluative
Kompetenz --- das schnelle Erkennen von Stilbrüchen, Faktenfehlern und
Tonalitätsabweichungen --- zur Kernkompetenz wird. Dies stellt keine
Dequalifizierung dar, sondern eine Neukalibrierung: Genre-Expertise und
Faktensicherheit bleiben unabdingbar, da ohne sie eine valide Bewertung
der KI-Entwürfe nicht möglich ist. Die publizistische Verantwortung ---
welcher Eintrag veröffentlicht wird --- verbleibt vollständig beim
Menschen.

\subsection{7.5.2 Skalierbarkeit und
White-Label}\label{skalierbarkeit-und-white-label}

Die White-Label-Architektur (\texttt{ef\_whitelabel}
vs.~\texttt{generic}) adressiert den in Kapitel 2.3 beschriebenen
strukturellen Wandel der Vereine zu eigenständigen Medienproduzenten.
Ein einzelnes System kann --- durch Instanzkonfiguration, Stilprofile
und Few-Shot-Referenzen --- verschiedene redaktionelle Stimmen bedienen,
ohne separate Codebases zu erfordern. Für Vereine mit begrenzten
Redaktionsressourcen senkt dies die Einstiegshürde in eine
professionelle Liveticker-Berichterstattung.

Die Zukunftsperspektive wird im Experteninterview (vgl.
Interviewleitfaden, Block 4, Fragen IF16--IF18) adressiert. Die
Skalierbarkeits- und White-Label-Bewertung stützt sich auf die in
Kapitel 5 dokumentierte Systemarchitektur und die Deployment-Erfahrung
mit der bestehenden \texttt{ef\_whitelabel}-Instanz; eine externe
Nutzerstudie ist in Abschnitt 8.3.2 vorgesehen.

\subsection{7.5.3 Ethische Überlegungen}\label{ethische-uxfcberlegungen}

Die automatisierte Generierung journalistischer Texte wirft die Frage
der Transparenz auf: Sollten Leser wissen, ob ein Ticker-Eintrag von
einem Menschen oder einer KI verfasst wurde? Das System speichert die
Herkunft jedes Eintrags im Feld \texttt{source} (\texttt{"ai"},
\texttt{"manual"}), legt die Kennzeichnung gegenüber dem Endnutzer aber
nicht offen. Für einen produktiven Einsatz wäre eine
Kennzeichnungspflicht --- etwa durch ein dezentes Label wie
„KI-unterstützt'' --- zu diskutieren, insbesondere im Kontext der
EU-KI-Verordnung (AI Act, EU 2024/1689), deren Transparenzpflichten für
KI-generierte Texte (Art. 50 Abs. 4) gemäß Art. 113 ab August 2026
gelten.

\begin{center}\rule{0.5\linewidth}{0.5pt}\end{center}

\section{Synthese}\label{synthese}

Die Diskussion der Evaluationsergebnisse entlang der vier Dimensionen
--- Abgrenzung zum Stand der Technik, methodische und technische
Limitationen, Betriebsmodi und Prompt-Architektur --- ergibt ein
konsistentes Gesamtbild: Das System löst die in Kapitel 1.1
identifizierten drei Problemdimensionen (Zeitdruck, Mehrsprachigkeit,
White-Label-Bedarf) durch technisch nachweisbare Mechanismen; die
verbleibenden Schwächen (Stil-Inkonsistenz, fehlende Authentifizierung,
eingeschränkte Stichprobengröße) sind klar lokalisiert und lösbar.

Die Dimension der Prompt-Architektur (Kapitel 7.4) verdeutlicht dabei,
dass die Textqualität nicht primär vom Modell abhängt, sondern von der
Qualität der Kontextübergabe: Few-Shot-Referenzen konditionieren Format
und Tonalität zuverlässig, während Score-Halluzinationen durch präzisere
Kontextstrukturierung weiter reduziert werden können.

Die zentrale Designentscheidung --- der \texttt{coop}-Modus als primärer
Betriebsmodus --- erweist sich aus dieser Diskussion heraus als
strukturell richtig: Er balanciert die Effizienzgewinne der KI mit der
publizistischen Verantwortung des Redakteurs und adressiert damit den in
Kapitel 1.2 formulierten Zielkonflikt zwischen Geschwindigkeit und
Qualität. Die Verschiebung der Redakteursrolle von der Textproduktion
zur Textkuration ist keine Umgehung journalistischer Sorgfalt, sondern
ihre Neukalibrierung unter den Bedingungen der
Echtzeit-Berichterstattung.

Die offenen Punkte --- externe Nutzerstudie, Authentifizierungsschicht,
A/B-Vergleich Few-Shot vs.~ohne --- definieren den nächsten
Entwicklungsschritt, der in Kapitel 8.3 als konkreter Ausblick
formuliert wird.
